% !TeX program = xelatex
% Madrid and 4:3 are IMPE's restored global Beamer defaults.
\documentclass{impebeamer}
\UseTemplateSet{features={tables,image}}
\title{Comparing Manuscript Witnesses}
\subtitle{From Transcription to Editorial Choice}
\author{Author Name}
\date{\today}
\begin{document}
\begin{frame}
  \titlepage
\end{frame}

\section{Question}
\begin{frame}{Roadmap}
  \tableofcontents
  \medskip
  \begin{block}{Purpose}
    A source comparison, methodological discussion and cautious
    editorial conclusion using entirely fictional witnesses.
  \end{block}
\end{frame}

\begin{frame}{Research Question}
  \begin{itemize}
    \item How should an editor represent competing readings?
    \item Which observations count as manuscript evidence?
    \item When is a preferred reading justified?
  \end{itemize}
  \begin{alertblock}{Guiding distinction}
    What survives in a source is not identical to what
    an editor reconstructs.
  \end{alertblock}
\end{frame}

\section{Evidence}
\begin{frame}{The Imagined Corpus}
  \begin{description}
    \item[Witness A] Earlier copy with a concise formulation.
    \item[Witness B] Later copy preserving an extra modifier.
    \item[Witness C] Fragmentary copy, damaged at the variant.
  \end{description}
  \medskip
  These letters do not imply a rank of reliability.
  Every witness needs its own material description.
\end{frame}

\begin{frame}{Reading the Evidence}
  A useful transcription distinguishes:
  \begin{enumerate}
    \item visible letters and punctuation;
    \item gaps, corrections and additions;
    \item editorial supplements and normalization.
  \end{enumerate}
  \medskip
  \begin{block}{Principle}
    Conjecture belongs in the apparatus, not silently
    in the description of surviving traces.
  \end{block}
\end{frame}

\begin{frame}{Three Readings at a Glance}
  \small
  \begin{table}
    \centering
    \begin{tabular}{lll}
      \toprule
      Source & Imagined text & Status \\
      \midrule
      A & the record is kept & complete \\
      B & the record is carefully kept & expanded \\
      C & the record \ldots{} kept & damaged \\
      \bottomrule
    \end{tabular}
  \end{table}
  \medskip
  C's missing middle cannot provide positive support
  for the shorter text transmitted by A.
\end{frame}

\section{Interpretation}
\begin{frame}{Two Possible Histories}
  \begin{columns}[T]
    \begin{column}{.47\textwidth}
      \begin{block}{Expansion}
        B may add a modifier to clarify
        A's abrupt instruction.
      \end{block}
    \end{column}
    \begin{column}{.47\textwidth}
      \begin{block}{Omission}
        A may have lost a modifier
        still preserved by B.
      \end{block}
    \end{column}
  \end{columns}
  \medskip
  Neither explanation follows from chronology alone.
\end{frame}

\begin{frame}{Testing an Editorial Hypothesis}
  \begin{itemize}
    \item Compare independent parallels and local vocabulary.
    \item Examine characteristic errors in each witness.
    \item Identify evidence that could overturn the inference.
  \end{itemize}
  \begin{exampleblock}{A revisable hypothesis}
    An independent witness with B's wording
    would strengthen the longer reading.
  \end{exampleblock}
\end{frame}

\begin{frame}{Presenting the Result}
  For a short reading edition we provisionally print A's form.
  \begin{block}{The apparatus records}
    \begin{itemize}
      \item B's additional modifier;
      \item C's damaged middle;
      \item the reasons for the adopted text.
    \end{itemize}
  \end{block}
  A synoptic edition may be better when variation
  itself is the research question.
\end{frame}

\section{Conclusion}
\begin{frame}{Takeaways}
  \begin{enumerate}
    \item Describe witnesses before selecting a text.
    \item Separate observed evidence from judgement.
    \item Make alternative readings recoverable.
  \end{enumerate}
  \medskip
  \centering
  \textbf{A readable edition should also be an inspectable one.}
\end{frame}
\end{document}
